% Back matter: Recommended Reading
% An author's reading list. Not a bibliography. A map for the
% reader who wants to go deeper into the territory the book crosses.

\chapter*{Recommended Reading}
\addcontentsline{toc}{chapter}{Recommended Reading}

These are the books and essays that shaped \emph{Demons at Work}, either as direct influences or as works that illuminate the territory the novella crosses. None of them are required. All of them are rewarding.

\bigskip

\textbf{George Eliot, \emph{Middlemarch} (1871--72).}
The novel Gordon reads aloud to an empty room. Eliot's subject is ordinary people making choices that feel small and turn out to be the architecture of their entire lives. Her argument is that the tragedy of most human existence is not dramatic but frequent: it lies in the gap between what people intend and what they accomplish. If this novella has a heart, it is borrowed from Eliot.

\bigskip

\textbf{Albert Camus, \emph{The Myth of Sisyphus} (1942).}
``One must imagine Sisyphus happy.'' Camus asks whether life is worth living when it has no transcendent meaning, and answers: yes, but not because it has meaning. Because you roll the boulder. The question of why Tert continues to do his work, and whether Gordon's weekly reading is the same act in a different register, is a question Camus asked first.

\bigskip

\textbf{Dante Alighieri, \emph{Inferno} (c.\,1314).}
The original hell-as-bureaucracy. Dante gave hell organizational structure, jurisdictional boundaries, assigned roles, and precise operational logic centuries before anyone put a demon in a cubicle. The \emph{Inferno} established that hell is not chaos. It is administration.

\bigskip

\textbf{C.\,S. Lewis, \emph{The Screwtape Letters} (1942).}
A senior demon writes letters of professional advice to a junior demon assigned to corrupt a human soul. Lewis understood that evil, seen from the inside, looks like office politics. \emph{Screwtape} is the direct genre predecessor to this novella, and the distance between the two books is the distance between satire that believes in the cosmic war and satire that does not know whether there is one.

\bigskip

\textbf{Herman Melville, ``Bartleby, the Scrivener'' (1853).}
``I would prefer not to.'' Bartleby is the worker who stops. He does not rebel, does not explain, does not leave. He simply ceases to participate and remains in the office. Tert is the opposite: the worker who sees everything Bartleby sees and keeps working. Whether this makes Tert more complicit or more human than Bartleby is a question both stories leave open.

\bigskip

\textbf{David Graeber, \emph{Bullshit Jobs: A Theory} (2018).}
An anthropologist's argument that a significant percentage of modern work is perceived as meaningless by the people who do it, and that this meaninglessness causes genuine psychological harm. Graeber's taxonomy of pointless work (flunkies, goons, duct tapers, box tickers, taskmasters) maps uncomfortably well onto hell's departmental structure.

\bigskip

\textbf{Thomas Nagel, ``What Is It Like to Be a Bat?'' (1974).}
Nagel's essay argues that consciousness has a subjective character that cannot be captured by objective description. You cannot know what it is like to be a bat by studying bat neurology. The horror genre does not ask what it is like to be the demon. This novella does.

\bigskip

\textbf{Franz Kafka, \emph{The Trial} (1925) and \emph{The Castle} (1926).}
Kafka's protagonists navigate bureaucratic systems whose rules are real, whose authorities are inaccessible, and whose purposes are never explained. The system is not broken. The system is working. The people inside it cannot tell the difference. Tert would recognize Josef K.'s predicament, though he would have opinions about K.'s approach.

\bigskip

\textbf{Kazuo Ishiguro, \emph{The Remains of the Day} (1989).}
A butler reflects on a career spent in devoted service to a man whose values were wrong. Stevens cannot separate professional excellence from moral complicity. His craft (the perfect household, the perfectly timed silver polishing) was the vehicle through which someone else's harm was made frictionless. Of all the books on this list, this may be the closest to what \emph{Demons at Work} is trying to do.

\bigskip

\textbf{John Gardner, \emph{Grendel} (1971).}
\emph{Beowulf} retold from the monster's perspective. Grendel watches humans build meaning through poetry, community, and narrative, and cannot participate. He is an existentialist trapped in someone else's epic. The tonal predecessor to Tert's voice: intelligent, isolated, watching a world that has not made room for him to be a person.

\bigskip

\textbf{Hannah Arendt, \emph{Eichmann in Jerusalem} (1963).}
The banality of evil. Arendt's argument is that enormous harm is administered not by monsters but by ordinary functionaries doing ordinary jobs within a system that has made compliance the path of least resistance. Tert is not Eichmann. But the structural question is the same: what does complicity look like when the work is routine, the forms are standard, and nobody in the building can explain the mechanism by which the harm occurs?

\bigskip

\textbf{Mary Shelley, \emph{Frankenstein} (1818).}
The creature's chapters. Shelley gives the monster a voice and the voice is articulate and the articulation makes everything harder, not easier. The creature can describe what it feels and what it wants, and the description changes nothing about how the world responds. Tert's narration operates under the same constraint: the interiority is precise, and the system has no use for it.

\bigskip

\textbf{Joan Didion, \emph{The Year of Magical Thinking} (2005).}
Grief rendered with the cold precision of a professional writer examining her own loss. Didion does not sentimentalize. She catalogs. The specificity of her observation (the shoes she cannot give away, the medical records she rereads) is the same specificity with which Tert catalogs Gordon's house, and with which Gordon continues to set a table for one.
